\chapter{States and sewing}
\label{mht:ch:states-sewing}

The residual action determines equations of motion and their gauge
symmetries. To obtain an amplitude, one must also choose a state.
For a boundary circle this requires a real form, a polarization,
a measure and a Hamiltonian. These choices determine the disk
amplitude and the pairing used when two boundaries are sewn.

We first prepare the compact residual gauge theory and compute its
character amplitudes. Holonomy density and invariant ordering then
exhibit two quantities that a normalized Gaussian vacuum does not
determine. Finite Gaussian integration, a free gauge quartet and an
infinite family of scalar oscillators provide explicit sewing maps.
A nonnegative interaction preserves scalar sewing at fixed positive
smoothing. Its divergent coupling series has a Borel reconstruction,
and a Hermite unitary identifies the interacting Fock Hamiltonian.
The final comparison asks which of these data the native Neumann
quantum theory supplies.

\section{The residual boundary and its real quantization}

Write $b\in\Omega^1(U,\mathfrak{gl}_N)$ and
$J\in\Omega^0(U,\mathfrak{gl}_N)$ for the physical residual fields.
The action of Proposition~\ref{mht:prop:exact-slice-action} is
\[
 S_{\mathrm{res}}=\kappa\int_U\operatorname{tr}(JF_b)
       -\frac\kappa2\int_{\partial U}\operatorname{tr}(bJ),
 \qquad F_b=d_Ub+b\wedge b.
\]
Its boundary one-form is
\[
 \Theta_{\partial U}=\frac\kappa2\int_{\partial U}
                    \operatorname{tr}(J\delta b-b\delta J).
\]
Adding $\kappa\int_{\partial U}\operatorname{tr}(bJ)/2$ changes
this primitive to $\kappa\int_{\partial U}\operatorname{tr}(J\delta b)$.
The boundary two-form is unchanged. We use the resulting connection
polarization, in which $b$ is the boundary configuration variable.

Choose a connected compact Lie group $G$ with Lie algebra
$\mathfrak g$, a positive invariant inner product
$\langle-,-\rangle$, and numerical parameters
$\kappa\in\mathbb R\setminus\{0\}$, $v>0$, $\hbar>0$.
Let $U$ be an oriented surface with a positive area form.
We declare the real Euclidean preparation by the first-order action
\begin{equation}\label{mht:eq:prepared-residual-action}
 S_E=\int_U\left(i\kappa\langle J,F_b\rangle
              +\frac v2\langle J,J\rangle\,d\mathrm{vol}\right).
\end{equation}
Here $J$ is a real adjoint scalar and $b$ is a $G$-connection.
The Gaussian measure in $J$ is normalized as part of this preparation.
Completing its square gives the density
$\kappa^2|F_b|^2/(2v)$ for the second-order action.

The corresponding canonical momentum is $P=\kappa J$.
Its kinetic energy per unit spatial length is
$v\langle P,P\rangle/(2\kappa^2)$.
Choose the positive Laplace--Beltrami quantization and assign energy
zero to the trivial representation. For the positive Laplacian
$\Delta_G$ of the chosen metric, this gives
\begin{equation}\label{mht:eq:prepared-heat-time}
 \begin{aligned}
 \mathsf H=\frac{v\hbar^2}{2\kappa^2}\Delta_G,
 \qquad &t=\frac{v\hbar}{\kappa^2}\operatorname{Area}(U),\\
 &e^{-\operatorname{Area}(U)\mathsf H/\hbar}=e^{-t\Delta_G/2}.
 \end{aligned}
\end{equation}

An additive energy constant would multiply each amplitude by an
area factor. The zero-energy convention fixes this freedom.
We also fix normalized Haar measure and the disk normalization below.
These are the representation-theoretic choices in two-dimensional
Yang--Mills theory.\footnote{E.~Witten, \emph{Two Dimensional Gauge
Theories Revisited}, \S4, including \S\S4.1--4.2,
\url{https://arxiv.org/abs/hep-th/9204083}. The Haar and vacuum
normalizations used here are specified explicitly.}

\subsection{The conjugation constraint}

Let $\mathcal R(G)$ be the complex representative functions on $G$.
Choose a basis of $\mathfrak g$ with structure constants
$f^a{}_{bc}$. Let $G_a$ denote the infinitesimal conjugation
operators, with $[G_a,G_b]=f^c{}_{ab}G_c$.
On the graded algebra
$\mathcal R(G)\otimes\Lambda(\mathfrak g^*)$, put
$|c^a|=1$ and use left odd derivatives. Its BRST differential is
\begin{equation}\label{mht:eq:residual-group-brst}
 Q=c^aG_a-\frac12f^a{}_{bc}c^bc^c\partial_{c^a}.
\end{equation}

\begin{proposition}\label{mht:prop:residual-group-brst}
The differential $Q$ squares to zero and commutes with
$e^{-t\Delta_G/2}$. Its degree-zero cohomology is
$\mathcal R(G)^G$. For normalized Haar measure $dg$, the Hilbert
completion of this cohomology is
\[
 \mathcal H_G=L^2(G,dg)^G.
\]
The orthogonal projection onto $\mathcal H_G$ is
$P_{\mathrm{inv}}f(U)=\int_Gf(hUh^{-1})\,dh$.
\end{proposition}

\begin{proof}
The coefficient of two ghosts in $Q^2$ is
$c^ac^b([G_a,G_b]-f^c{}_{ab}G_c)/2$.
The coefficient of three ghosts is the Jacobi identity.
Both vanish on the whole exterior algebra.
The bi-invariant Laplacian commutes with each $G_a$ and acts
trivially on the ghosts, proving the heat commutation relation.

In degree zero, $Qf=0$ is the system $G_af=0$ for all $a$.
Connectedness of $G$ identifies its solutions with conjugation
invariants. There are no boundaries from negative degrees.
Haar invariance proves that $P_{\mathrm{inv}}$ is idempotent
and self-adjoint. Jensen's inequality gives its $L^2$ bound.
Peter--Weyl approximation, followed by this projection, shows
that invariant representative functions are dense in its range.
\end{proof}

This identifies ghost degree zero. Higher Lie algebra cohomology
is retained by the BRST complex and need not vanish.

\subsection{Disk, cylinder and surface amplitudes}

For $\lambda\in\widehat G$, let $d_\lambda$, $\chi_\lambda$
and $C_2(\lambda)$ denote its dimension, character and positive
Laplacian eigenvalue. The characters are an orthonormal basis
of $\mathcal H_G$. Equip this space on its smooth core with
Haar convolution,
\[
 (f*g)(U)=\int_G f(UW^{-1})g(W)\,dW.
\]
The preparation uses the complete sum over $\widehat G$.
It has no additional discrete theta weight or fixed-bundle
projection. For nonsimply connected $G$, prescribing such a
sector would change this global choice.

\begin{theorem}\label{mht:thm:residual-character-sewing}
At positive heat time, the cylinder kernel and disk state are
\begin{align}
 K_t(U,V)&=\sum_{\lambda\in\widehat G}
 e^{-tC_2(\lambda)/2}\chi_\lambda(U)\overline{\chi_\lambda(V)},
       \label{mht:eq:residual-cylinder}\\
 D_t(U)&=\sum_{\lambda\in\widehat G}
 d_\lambda e^{-tC_2(\lambda)/2}\chi_\lambda(U).
       \label{mht:eq:residual-disk}
\end{align}
They are smooth, and for $s,t>0$ satisfy
\begin{equation}\label{mht:eq:residual-cylinder-sewing}
 \int_GK_s(U,W)K_t(W,V)\,dW=K_{s+t}(U,V).
\end{equation}
The disk tends to the identity point source distributionally
as $t\downarrow0$. With this normalization, the amplitude
of a genus-$g$ surface with $b$ outgoing boundaries is
\begin{equation}\label{mht:eq:residual-surface-amplitude}
 Z_{g,b}(t;U_1,\ldots,U_b)
 =\sum_\lambda d_\lambda^{2-2g-b}e^{-tC_2(\lambda)/2}
                      \prod_{i=1}^b\chi_\lambda(U_i).
\end{equation}
An incoming boundary carries the conjugate character.
These amplitudes satisfy boundary and handle sewing.
\end{theorem}

\begin{proof}
The heat operator is diagonal in the character basis, giving
\eqref{mht:eq:residual-cylinder}. The identity point source has
Fourier coefficients $\chi_\lambda(1)=d_\lambda$.
Heat evolution therefore gives \eqref{mht:eq:residual-disk}.
Character orthogonality proves
\eqref{mht:eq:residual-cylinder-sewing} for finite sums.

Highest weights have lattice growth. Representation dimensions
and the coefficients of differentiated matrix elements grow
polynomially in the weight. The Casimir grows quadratically
for a positive invariant metric. Consequently the sums and
all their derivatives converge locally uniformly for $t>0$.
These bounds justify Haar integration and prove smoothness.
On a smooth test function the Fourier coefficients also give
the asserted distributional limit of the disk.

Schur orthogonality gives
\[
 \chi_\lambda*\chi_\mu
       =\delta_{\lambda\mu}\frac{\chi_\lambda}{d_\lambda}.
\]
Thus a three-holed sphere contributes $d_\lambda^{-1}$.
Its product with the disk coefficient $d_\lambda$ is the
unit cylinder coefficient. Decomposition into these pieces
gives \eqref{mht:eq:residual-surface-amplitude}.

Sewing opposed boundary characters identifies their labels.
For two surfaces, their dimension exponents add to the Euler
characteristic of the glued surface. Gluing two boundaries of
one surface leaves that exponent unchanged: $g$ increases by
one and $b$ decreases by two. These observations prove both
sewing laws. The positive-time bounds justify each passage
from finite sums to the displayed amplitudes.
\end{proof}

In particular, the sphere and torus amplitudes are
\begin{equation}\label{mht:eq:residual-closed-amplitudes}
 Z_{0,0}(t)=\sum_\lambda d_\lambda^2e^{-tC_2(\lambda)/2},
 \qquad
 Z_{1,0}(t)=\operatorname{Tr}_{\mathcal H_G}e^{-t\Delta_G/2}
           =\sum_\lambda e^{-tC_2(\lambda)/2}.
\end{equation}
For a nontrivial connected compact group, the torus trace is
strictly greater than one. If $\dim G>0$, the zero-area disk
is not an $L^2$ state. Indeed, invariant smooth functions equal
to one at the identity can have arbitrarily small $L^2$ norm,
by shrinking an invariant bump function to the identity.
Point evaluation is therefore unbounded on the invariant core.
For the trivial group the point source is the normalized constant.

On a closed surface the positive summands define a probability
distribution after division by $Z_{g,0}(t)$. Differentiation gives
\begin{equation}\label{mht:eq:residual-thermal-response}
 -2\partial_t\log Z_{g,0}=\langle C_2\rangle,
 \qquad
 4\partial_t^2\log Z_{g,0}
       =\langle C_2^2\rangle-\langle C_2\rangle^2.
\end{equation}
The same positive-time majorants justify both derivatives.
These are moments of the chosen quadratic quantum operator.
Higher invariant operators require an ordering convention.

\subsection{A normalization and tail bound for \texorpdfstring{$SU(2)$}{SU(2)}}

Give $SU(2)$ its round-$S^3$ metric. Its irreducibles have labels
$n\geq0$ and data
\[
 d_n=n+1,\qquad C_2(n)=n(n+2),\qquad
 \chi_n(\theta)=\frac{\sin((n+1)\theta)}{\sin\theta}.
\]
The endpoint values of the quotient are its continuous limits.
On class functions, normalized Haar measure is
$(2/\pi)\sin^2\theta\,d\theta$ for $0\leq\theta\leq\pi$.
The character inner product is therefore the elementary sine
inner product.

At $t=1$, the convergent sums give
\[
 Z_{1,0}(1)\simeq1.24200505276743,
 \qquad Z_{0,0}(1)\simeq2.06636525163437.
\]
The corresponding torus mean and variance of $C_2$ are
approximately $0.66373337$ and $2.22320420$.

To control truncation, set $\alpha=t/2$ and retain $n\leq K$,
where $K\geq0$ is integral. Write $\operatorname{erfc}$ for the
complementary error function. Since $n(n+2)\geq n^2$,
\begin{equation}\label{mht:eq:residual-torus-tail}
 0\leq\sum_{n>K}e^{-tn(n+2)/2}
 \leq\int_K^\infty e^{-\alpha x^2}\,dx
 =\frac{\sqrt\pi}{2\sqrt\alpha}
                         \operatorname{erfc}(\sqrt\alpha K).
\end{equation}

For $K+1\geq\alpha^{-1/2}$, the function
$x^2e^{-\alpha x^2}$ decreases on $[K+1,\infty)$.
Putting $m=n+1$ and integrating gives
\begin{align}
 \sum_{n>K}(n+1)^2e^{-tn(n+2)/2}
 &\leq e^\alpha\int_{K+1}^\infty x^2e^{-\alpha x^2}\,dx
       \notag\\
 &=e^\alpha\left[
 \frac{(K+1)e^{-\alpha(K+1)^2}}{2\alpha}
 +\frac{\sqrt\pi\operatorname{erfc}(\sqrt\alpha(K+1))}
                                     {4\alpha^{3/2}}\right].
       \label{mht:eq:residual-sphere-tail}
\end{align}
Because $|\chi_n|\leq n+1$, the same bound controls the
absolute disk-character tail. Both estimates require $t>0$.

\section{Holonomy density and invariant ordering}

The heat state uses a group holonomy and a quantum Hamiltonian.
Exponential coordinates contribute a Jacobian, while ordering
changes products of invariant insertions. We compute both with
their normalizations before asking for a map from the polar fields.

\subsection{The exponential half-density}

Let $\mathfrak g$ now be a finite-dimensional unimodular complex
Lie algebra. For a dimensionless exponential coordinate $X$ near
zero, write $L=\operatorname{ad}_X$. Left trivialization gives
\[
 \operatorname{dexp}_X=\int_0^1e^{-tL}\,dt
          =e^{-L/2}\frac{\sinh(L/2)}{L/2}.
\]
The quotient denotes its entire power series. Normalize the
tangent and invariant volume forms to agree at the identity.
Choose every square-root branch to have value one there.

\begin{proposition}\label{mht:prop:holonomy-half-density}
The exponential Jacobian and its half-density are
\begin{equation}\label{mht:eq:holonomy-half-density}
 j(X)=\det_{\mathfrak g}\frac{\sinh(L/2)}{L/2},\qquad
 \mathfrak j(X)=\exp\!\left(
 \frac12\operatorname{Tr}_{\mathfrak g}
                         \log\frac{\sinh(L/2)}{L/2}\right).
\end{equation}
On a sufficiently small neighborhood of zero,
\begin{equation}\label{mht:eq:holonomy-mode-product}
 \mathfrak j(X)=\lim_{K\to\infty}
 \prod_{n=1}^K\det_{\mathfrak g}
       \left(1+\frac{L^2}{4\pi^2n^2}\right)^{1/2}
\end{equation}
locally uniformly. Let $\mathfrak j_K$ denote this product through $n=K$.
If $L$ is real diagonalizable with real
eigenvalues, then, for $K\geq1$,
\begin{equation}\label{mht:eq:holonomy-mode-tail}
 0\leq\log\mathfrak j-\log\mathfrak j_K
       \leq\frac{\operatorname{Tr}_{\mathfrak g}L^2}{8\pi^2K}.
\end{equation}
\end{proposition}

\begin{proof}
Unimodularity gives $\operatorname{Tr}L=0$, so
$\det(e^{-L/2})=1$. This proves the Jacobian formula.
The scalar product for $\sinh(z/2)/(z/2)$ gives
\eqref{mht:eq:holonomy-mode-product} on eigenvalues, including
their algebraic multiplicities. Each logarithmic factor is
locally bounded by a constant times $n^{-2}$. Hence the
product converges locally uniformly on its normalized branch.

For real eigenvalues $\lambda_a$, apply
$0\leq\log(1+u)\leq u$ to
$u=\lambda_a^2/(4\pi^2n^2)$. The logarithmic tail is bounded by
$\sum_a\lambda_a^2\sum_{n>K}n^{-2}/(8\pi^2)$.
The estimate $\sum_{n>K}n^{-2}\leq K^{-1}$ proves the result.
\end{proof}

On a compact real form, the eigenvalues of $L$ are imaginary
and the factors become sine factors. The branch is local,
before their first zeros. For a quadratic Lie algebra, the
same half-determinant occurs in the effective action of
one-dimensional Chern--Simons theory.\footnote{A.~Alekseev
and P.~Mnev, \emph{One-dimensional Chern--Simons theory},
Proposition~1, \url{https://arxiv.org/abs/1005.2111}.}
That formula concerns the circle field complex. A comparison
with the Neumann sphere must identify the corresponding fields.

Expansion at $X=0$ gives
\begin{equation}\label{mht:eq:holonomy-density-expansion}
 \log\mathfrak j(X)=\frac1{48}\operatorname{Tr}L^2
       -\frac1{5760}\operatorname{Tr}L^4
       +\frac1{362880}\operatorname{Tr}L^6+O(X^8).
\end{equation}
For $\mathfrak g=\mathfrak{gl}_N$, put
$p_r=\operatorname{tr}(X^r)$, with $p_0=N$.

\begin{proposition}\label{mht:prop:holonomy-block-defect}
For every integer $m\geq0$,
\begin{equation}\label{mht:eq:adjoint-trace-polynomial}
 \operatorname{Tr}_{\mathfrak{gl}_N}(\operatorname{ad}_X)^m
       =\sum_{r=0}^m(-1)^r\binom mr p_{m-r}p_r.
\end{equation}
In particular, $\operatorname{Tr}L^2=2Np_2-2p_1^2$ and
$\operatorname{Tr}L^4=2Np_4-8p_1p_3+6p_2^2$.
For matrices $X,Y$ of sizes $N,M$, the quadratic part of
$\log\mathfrak j(X\oplus Y)-\log\mathfrak j(X)
-\log\mathfrak j(Y)$ is
\begin{equation}\label{mht:eq:holonomy-block-defect}
 \frac1{24}\left(M\operatorname{tr}X^2
        +N\operatorname{tr}Y^2
        -2\operatorname{tr}X\operatorname{tr}Y\right).
\end{equation}
\end{proposition}

\begin{proof}
Left and right multiplication on $\operatorname{End}(\mathbb C^N)$
commute. Their difference is $\operatorname{ad}_X$, and
$\operatorname{Tr}(L_X^aR_X^b)
=\operatorname{tr}(X^a)\operatorname{tr}(X^b)$.
The binomial theorem proves \eqref{mht:eq:adjoint-trace-polynomial}
without diagonalizing $X$.
Substituting $m=2,4$ gives the displayed traces.
The quadratic half-density logarithm is $(Np_2-p_1^2)/24$.
Its block-sum difference is \eqref{mht:eq:holonomy-block-defect}.
\end{proof}

The two off-diagonal Hom spaces between the blocks contribute
to this difference. Thus this quantum density is not an additive
family of classical invariant polynomials. Its holonomy coordinate
$X$ has not been identified with the residual matrix $J$.

\subsection{The universal quadratic and quartic contacts}

Use Hermitian generators in the complex Rees enveloping algebra,
completed in $\hbar$, with relations
\[
 [\widehat J_i,\widehat J_j]
       =i\hbar\epsilon_{ijk}\widehat J_k.
\]
Put $p=x_1^2+x_2^2+x_3^2$, $C=\sum_i\widehat J_i^2$, and
$\Delta=\sum_i\partial_{x_i}^2$.
PBW symmetrization $\operatorname{sym}$ averages all orderings
of a commutative monomial in the generators $\widehat J_i$.

In this convention define
\begin{equation}\label{mht:eq:su2-duflo-factor}
 \operatorname{Duf}_\hbar
 =\operatorname{sym}\circ
 \left(1+\frac{\hbar^2}{24}\Delta
          +\frac{\hbar^4}{1920}\Delta^2+\cdots\right).
\end{equation}
The series in parentheses is
$\sinh(\hbar\sqrt\Delta/2)/(\hbar\sqrt\Delta/2)$,
interpreted as a power series in $\Delta$.
The adjoint eigenvalues in the unscaled Hermitian basis are
$0,\sqrt p,-\sqrt p$, so this is the half-Jacobian factor.
The Duflo theorem identifies its invariant restriction as an
algebra isomorphism.\footnote{A.~Alekseev and E.~Meinrenken,
\emph{The non-commutative Weil algebra}, \S6.1,
\url{https://arxiv.org/abs/math/9903052}.}

\begin{proposition}\label{mht:prop:su2-ordering-contact}
The universal Rees algebra satisfies
\begin{align}
 \operatorname{Duf}_\hbar(p)&=C+\hbar^2/4,
       \label{mht:eq:su2-quadratic-ordering}\\
 \operatorname{sym}(p^2)&=C^2-\hbar^2C/3,
       \label{mht:eq:su2-quartic-contact}\\
 \operatorname{Duf}_\hbar(p^2)&=(C+\hbar^2/4)^2.
       \label{mht:eq:su2-duflo-product}
\end{align}
For the multiplicative ordering $p\mapsto C$, the collapsed
double-insertion contact relative to raw PBW symmetrization
is $\hbar^2C/3$.
\end{proposition}

\begin{proof}
The relations make $C$ central and give
\[
 \sum_j[\widehat J_j,[\widehat J_j,\widehat J_i]]
       =2\hbar^2\widehat J_i,
 \qquad
 \sum_j\widehat J_j\widehat J_i\widehat J_j
       =(C-\hbar^2)\widehat J_i.
\]
Grouping the six orders of two repeated indices yields
\[
 \operatorname{sym}(p^2)=\frac13\sum_{i,j}
 \left(\widehat J_i^2\widehat J_j^2
       +\widehat J_i\widehat J_j\widehat J_i\widehat J_j
       +\widehat J_i\widehat J_j^2\widehat J_i\right).
\]
The three sums equal $C^2$, $C^2-\hbar^2C$ and $C^2$.
This proves \eqref{mht:eq:su2-quartic-contact} before any
representation is chosen. Finally,
$\Delta p=6$, $\Delta p^2=20p$, and $\Delta^2p^2=120$.
Substitution into \eqref{mht:eq:su2-duflo-factor} gives
the other two identities.
\end{proof}

The invariant ordering with zero vacuum energy is
\[
 \mathcal Q(F(p))=\operatorname{Duf}_\hbar(F(p-\hbar^2/4)),
 \qquad \mathcal Q(p)=C.
\]
Translation in $p$ and the Duflo theorem prove multiplicativity.
In spin $j$, the quadratic eigenvalues before and after the
Duflo shift are $\hbar^2j(j+1)$ and $\hbar^2(j+1/2)^2$.
At numerical $\hbar=1$, the two half-Casimir evolutions differ
by $e^{-t/8}$. This scalar factor respects addition of $t$,
so sewing does not select between the two vacuum energies.
The round-$S^3$ Casimir above is four times this spin Casimir.

\section{Gaussian states with their normalization lines}

The distinction between a state and its normalized correlations
already occurs in one variable. Gaussian integration gives a
nonzero quadratic moment, while its value on the unit carries
a square root of a determinant. Under elimination, the moment
and determinant compose by different parts of the same block
matrix identity.

\subsection{The polynomial Gaussian complex}

Let $R=\mathbb C[\hbar]$ and let $Q$ be an invertible symmetric
complex $n\times n$ matrix. Set
\[
 B_Q=R[x_1,\ldots,x_n]\otimes\Lambda(\xi_1,\ldots,\xi_n),
 \qquad |x_i|=0,\quad |\xi_i|=-1.
\]
Using left odd derivatives, define
\begin{equation}\label{mht:eq:gaussian-state-differential}
 D_Q=\sum_i\bigl(\hbar\partial_{x_i}-(Qx)_i\bigr)
                                                   \partial_{\xi_i}.
\end{equation}

The even factors commute because $Q$ is symmetric, whereas
the odd derivatives anticommute. Thus $D_Q^2=0$.
These formulas define $B_Q,D_Q$ for every symmetric matrix.
Invertibility is required for the integration map below.

Over the space of nondegenerate symmetric matrices, let
$\mathscr G_Q$ be the degree-zero Gaussian square-root line.
Its local sections $\gamma_Q=(\det Q)^{-1/2}$ have the sign
transitions of the square root. On real positive matrices use
the positive section, corresponding to the density
$dx/(2\pi\hbar)^{n/2}$ for numerical $\hbar>0$.
Continuation retains this line and its chosen branch.

\begin{proposition}\label{mht:prop:sourced-gaussian-state}
The $R$-linear map
\begin{equation}\label{mht:eq:sourced-gaussian-state}
 \Pi_Q(F)=\gamma_Q\left.
 \exp\!\left(\frac\hbar2
       \sum_{i,j}(Q^{-1})_{ij}\partial_{x_i}\partial_{x_j}\right)
                       F(x,0)\right|_{x=0}
\end{equation}
is a cochain map from $(B_Q,D_Q)$ to $\mathscr G_Q\otimes R$
with zero differential. For commuting formal sources $t_i$,
\begin{equation}\label{mht:eq:gaussian-state-ward}
 Z_Q(t)=\gamma_Qe^{\hbar t^{\mathsf T}Q^{-1}t/2},
 \qquad
 \sum_jQ_{ij}\partial_{t_j}Z_Q=\hbar t_iZ_Q.
\end{equation}
These identities hold coefficientwise in
$\mathscr G_Q\otimes R[[t_1,\ldots,t_n]]$.
\end{proposition}

\begin{proof}
Write $\mathbb E_Q=\gamma_Q^{-1}\Pi_Q$ in a local branch.
The exponential terminates on each polynomial.
Commuting multiplication by $x_j$ through it gives
\[
 \mathbb E_Q(x_jf)=\hbar\sum_l(Q^{-1})_{jl}
                                      \mathbb E_Q(\partial_{x_l}f).
\]

Consequently $\Pi_QD_Q(\xi_i f)=0$.
An input with at least two odd factors retains an odd factor
after $D_Q$ and maps to zero. An input with no odd factors
has zero differential. This proves the cochain assertion.

Applying the heat exponential to the formal source gives
the first identity in \eqref{mht:eq:gaussian-state-ward}.
Differentiation proves the second.
\end{proof}

The Ward equations allow every nonzero source-independent
multiple of $Z_Q$. Integration data fix $Z_Q(0)=\gamma_Q$.
Even after normalization, evaluation at zero is insufficient:
for $Q=1$ in one variable,
\[
 D_Q(x\xi)=\hbar-x^2,
 \qquad \mathbb E_Q(x^2)=\hbar.
\]

For a real symmetric invertible Hessian $H$, use the damped
matrix $Q_\epsilon=\epsilon I-iH$, $\epsilon>0$.
Continuation through $\operatorname{Re}Q_\epsilon>0$ gives
\begin{equation}\label{mht:eq:gaussian-state-fresnel}
 \gamma_{-iH}=|\det H|^{-1/2}
                       e^{i\pi\operatorname{sign}(H)/4},
 \qquad \mathbb E_{-iH}(x_ix_j)=i\hbar(H^{-1})_{ij}.
\end{equation}
Here $\operatorname{sign}(H)$ is the number of positive
eigenvalues minus the number of negative eigenvalues.
Diagonalization reduces both formulas to the one-dimensional
Gaussian. No infinite determinant occurs in these identities.

\subsection{The full residual complex under elimination}

Split the even coordinates into $(x,y)$ and their odd partners
into $(\xi,\eta)$. For a symmetric matrix write
\[
 \mathsf Q=\begin{pmatrix}A&B\\B^{\mathsf T}&C\end{pmatrix},
 \qquad C_{\mathrm{eff}}=C-B^{\mathsf T}A^{-1}B,
\]
assuming that $A$ is invertible. The even shear and its
cotangent lift are
\begin{equation}\label{mht:eq:gaussian-state-shear}
 x'=x+A^{-1}By,\quad y'=y,\quad
 \xi'=\xi,\quad \eta'=\eta-B^{\mathsf T}A^{-1}\xi.
\end{equation}
They preserve the quantum contraction
$\sum_i\partial_{x_i}\partial_{\xi_i}
+\sum_j\partial_{y_j}\partial_{\eta_j}$ and have Berezinian one.
The quadratic form becomes
$x'^{\mathsf T}Ax'+y^{\mathsf T}C_{\mathrm{eff}}y$.

\begin{proposition}\label{mht:prop:gaussian-state-sewing}
Apply $\Pi_A$ in the primed eliminated variables and retain
the full $(y,\eta')$ complex. The result is a cochain map
\[
 (B_{\mathsf Q},D_{\mathsf Q})\longrightarrow
 \mathscr G_A\otimes(B_{C_{\mathrm{eff}}},D_{C_{\mathrm{eff}}}).
\]
Here the polynomial differential defines the residual complex
even if $C_{\mathrm{eff}}$ is degenerate. On formal sources,
the map is
\begin{equation}\label{mht:eq:partial-gaussian-state-source}
 e^{t_x^{\mathsf T}x+t_y^{\mathsf T}y}\longmapsto
 \gamma_Ae^{\hbar t_x^{\mathsf T}A^{-1}t_x/2}
             e^{(t_y-B^{\mathsf T}A^{-1}t_x)^{\mathsf T}y}.
\end{equation}
When $C_{\mathrm{eff}}$ is invertible, successive integration
agrees with $\Pi_{\mathsf Q}$ under the branch identification
\[
 \mathscr G_{\mathsf Q}\simeq
                  \mathscr G_A\otimes\mathscr G_{C_{\mathrm{eff}}}.
\]
Iterated integration is associative on that branch.
\end{proposition}

\begin{proof}
The inverse-transpose transformation of the odd variables
splits $D_{\mathsf Q}$ into the $A$ and $C_{\mathrm{eff}}$
differentials. Proposition~\ref{mht:prop:sourced-gaussian-state}
therefore proves the cochain assertion with all residual odd
variables retained. Substituting $x=x'-A^{-1}By$ proves
\eqref{mht:eq:partial-gaussian-state-source}.

The block inverse formula identifies the source exponent
after the second integration. The identity
$\det\mathsf Q=\det A\det C_{\mathrm{eff}}$ identifies the
normalization lines. Fubini on the positive real domain fixes
their common branch and proves associativity.
Analytic continuation preserves these identities.
For real oscillatory Hessians, the same congruence adds
signatures, so the Fresnel phases compose as well.
\end{proof}

For $\mathsf Q=\left(\begin{smallmatrix}2&1\\1&3\end{smallmatrix}\right)$,
one has $C_{\mathrm{eff}}=5/2$ and
\[
 \frac\hbar{10}(3t_x^2-2t_xt_y+2t_y^2)
       =\frac\hbar4t_x^2+\frac\hbar5(t_y-t_x/2)^2.
\]
The vacuum factor is $5^{-1/2}=2^{-1/2}(5/2)^{-1/2}$.
A kernel of $C_{\mathrm{eff}}$ instead remains in the residual
complex. The partial integration never inverts that kernel.

\subsection{Continuous Wick preparation on regular kernels}

For infinitely many fields, specify a test space before applying
the same contraction. Let $(W,d)$ be a nuclear graded test
complex with a continuous graded-symmetric pairing $b$ of
degree one. Assume
\[
 b(du,v)+(-1)^{|u|}b(u,dv)=0.
\]
Let $(H,d_H)$ be finite-dimensional. Choose continuous maps
$i,p,h$ of degrees $0,0,-1$, with $i,p$ cochain maps, such that
\[
 pi=1,\qquad dh+hd=1-ip,\qquad h^2=hi=ph=0.
\]
Assume $P=ip$ is self-adjoint for $b$ and
$C(u,v)=-b(hu,v)$ is graded symmetric of degree zero.

Use the completed nuclear tensor spaces of the tests.
For a smooth test complex these are smooth compact kernels,
and covariance distributions act continuously on them.
Let $\operatorname{Sym}_{\mathrm{reg}}^nW$ be the symmetric
part of the chosen $n$-fold tensor space. Denote unordered
pair contraction with $C$ by $\partial_C$, and pair contraction
with $b$ by $\Delta_b$. Put $b_H(a,a')=b(ia,ia')$.

\begin{proposition}\label{mht:prop:continuous-wick-state}
The normalized Wick map
\[
 \Gamma_h=\operatorname{Sym}(p)e^{\hbar\partial_C}
\]
is a continuous cochain map
\[
 \left(\bigoplus_{n\geq0}\operatorname{Sym}_{\mathrm{reg}}^nW[\hbar],
                                      d+\hbar\Delta_b\right)
 \longrightarrow
 \left(\operatorname{Sym}H[\hbar],d_H+\hbar\Delta_{b_H}\right).
\]
It sends $1$ to $1$ and extends to coefficientwise
$\hbar$-series in these spaces.
\end{proposition}

\begin{proof}
On homogeneous generators,
\begin{align*}
 [d,\partial_C](uv)
   &=-C(du,v)-(-1)^{|u|}C(u,dv)\\
   &=b((dh+hd)u,v)=b(u,v)-b(Pu,Pv).
\end{align*}
The last equality uses self-adjointness and $P^2=P$.
Summing over pairs gives
$[d,\partial_C]=\Delta_b-\Delta_{b(P-,P-)}$.
Constant contractions supercommute, and $\partial_C$ is even.
Thus
\[
 e^{\hbar\partial_C}(d+\hbar\Delta_b)e^{-\hbar\partial_C}
          =d+\hbar\Delta_{b(P-,P-)}.
\]
Applying $\operatorname{Sym}(p)$ proves the cochain equation.

Each term contracts disjoint pairs of variables in one smooth
compact kernel. Partial distributional evaluation leaves a
smooth kernel with support in the projected compact set.
The nuclear tensor topology makes these operations continuous.
There is no multiplication of singular distributions on a
common diagonal. At fixed field number the exponential
terminates. At fixed $\hbar$ order only finitely many input
coefficients contribute, proving the formal extension.
\end{proof}

Strong convergence of the covariance distributions on each
fixed compact test space gives coefficientwise convergence
of these maps. Their value on the unit remains one.
An unnormalized state still needs its determinant section
and a boundary-compatible family of contractions.

\section{A free gauge state with nonzero contractions}

The preceding Gaussian cochain map is compatible with a quantum
second-order differential. A free BRST differential instead
acts as a derivation. A quartet gives an explicit preparation
for this case and separates cancellation of the graded trace
from the values of its observables.

Over $\mathbb C[[\hbar]]$, set
\[
 A=\mathbb C[x,b]\otimes\Lambda(\bar c,c)[[\hbar]],
 \qquad |x|=|b|=0,\quad |\bar c|=-1,\quad |c|=1.
\]
Each coefficient of $\hbar$ is polynomial in the displayed
variables. Use the odd order $\bar c,c$ and left derivatives.
Define
\[
 Q=c\partial_x+b\partial_{\bar c},\qquad
 s=x\partial_c+\bar c\partial_b.
\]

Let $N$ count total polynomial and exterior degree.
The identities $Q^2=s^2=0$ and $Qs+sQ=N$ hold on generators,
and hence on the algebra by the derivation rule.
Put $h_0=s/N$ on positive degree and zero on constants.
If $i$ includes constants, then
$Qh_0+h_0Q=1-i\operatorname{ev}_0$.

Introduce the even contraction
\begin{equation}\label{mht:eq:quartet-wick}
 C=\partial_x\partial_b+\partial_c\partial_{\bar c},
 \qquad W=e^{\hbar C},\qquad P_G=\operatorname{ev}_0W.
\end{equation}
The exponential terminates on each polynomial and is a
continuous linear automorphism of the underlying formal module.

\begin{theorem}\label{mht:thm:quartet-preparation}
The functional $P_G$ is $Q$-closed and has moments
\begin{align}
 P_G(x^mb^n)&=\delta_{mn}m!\hbar^m,\notag\\
 P_G(\bar c c\,x^mb^n)&=\hbar\delta_{mn}m!\hbar^m.
       \label{mht:eq:quartet-moments}
\end{align}
Terms with exactly one ghost have zero expectation.
For $H_G=W^{-1}h_0W$,
\begin{equation}\label{mht:eq:quartet-contraction}
 QH_G+H_GQ=1-iP_G,
 \qquad H_G^2=H_Gi=P_GH_G=0.
\end{equation}
\end{theorem}

\begin{proof}
The commutators of $Q$ with the two terms of $C$ are
$-\partial_x\partial_{\bar c}$ and
$\partial_x\partial_{\bar c}$. Hence $[Q,C]=0$.
Evaluation at zero kills $Q$, so $P_GQ=0$.

The bosonic contraction survives evaluation only when it
differentiates $x^mb^n$ equally often in both variables.
For $m=n$ its coefficient is $(m!)^2\hbar^m/m!$.
Also $\partial_c\partial_{\bar c}(\bar c c)=1$.
These commuting contractions prove the moment formulas.
Conjugation of the $h_0$ identities by $W$ gives
\eqref{mht:eq:quartet-contraction}, since $Wi=i$.
\end{proof}

For example,
\[
 Q(\bar c x)=bx-\bar c c,\qquad
 P_G(bx)=P_G(\bar c c)=\hbar.
\]
The Ward identity equates two nonzero contractions.
The preparation is a graded algebraic functional, with no
claim of a positive measure on odd variables.

\subsection{Transfer and its Hilbert completion}

For $0<r<1$, define $T_r=W^{-1}r^NW$ on polynomials.
Because $[Q,N]=[Q,W]=0$,
\begin{equation}\label{mht:eq:quartet-transfer}
 [Q,T_r]=0,\qquad T_{r_1}T_{r_2}=T_{r_1r_2},
 \qquad P_GT_r=P_G,\qquad T_r1=1.
\end{equation}

For an analytic trace, fix a numerical value of $\hbar$.
Declare the vectors
\[
 W^{-1}(x^ab^d\bar c^{\nu_1}c^{\nu_2}),
 \qquad a,d\geq0,\quad \nu_1,\nu_2\in\{0,1\},
\]
orthonormal, with parity $\nu_1+\nu_2$.
This defines the Hilbert completion. The operator $T_r$ is
diagonal there, with eigenvalue $r^{a+d+\nu_1+\nu_2}$, so
\begin{equation}\label{mht:eq:quartet-supertrace}
 \operatorname{Tr}|T_r|=\frac{(1+r)^2}{(1-r)^2},
 \qquad \operatorname{Str}T_r=1.
\end{equation}
The differential and homotopy identities hold on the invariant
polynomial core. No boundedness of $W$ on a different Hilbert
space is assumed.

This free differential also has a finite BFV charge.
Give the cotangent momenta of $(x,b,\bar c,c)$ complementary
degrees so that the canonical Poisson bracket has degree zero.
Let $\Theta$ be the degree-one fiber-linear Hamiltonian of $Q$.
The correspondence from vector fields to these Hamiltonians
preserves brackets, so $\{\Theta,\Theta\}=0$.
Coordinate quantization gives the square-zero operator $\hbar Q$.

\subsection{Nonzero Fourier modes and the residual sector}

Take a flat product
\[
 Y=S^1_s\times(\mathbb C/\Lambda)^2\cong\mathbb T^5,
 \qquad D=d_s+\bar\partial.
\]
For a real dual-lattice vector
$k=(k_s,k_{x_1},k_{y_1},k_{x_2},k_{y_2})$, the mode complex is
$\Lambda^\bullet\mathbb C^3$ with differential
$D_k=\eta(k)\wedge$, where
\[
 \eta(k)=ik_s\,ds
       +\frac{ik_{x_1}-k_{y_1}}2\,d\bar z_1
       +\frac{ik_{x_2}-k_{y_2}}2\,d\bar z_2.
\]

For $k\ne0$, this covector is nonzero. The Hermitian dual
vector $\eta(k)^\sharp$ satisfies
$\eta(k)(\eta(k)^\sharp)=\|\eta(k)\|^2$, and
\begin{equation}\label{mht:eq:mixed-fourier-contraction}
 h_k=\frac{\iota_{\eta(k)^\sharp}}{\|\eta(k)\|^2},
 \qquad D_kh_k+h_kD_k=1.
\end{equation}
To see the exact pairs, choose
$u_1=\eta(k)/\|\eta(k)\|$ and an orthonormal complement
$u_2,u_3$. For each $I\subset\{2,3\}$,
\[
 D_ku_I=\|\eta(k)\|u_1\wedge u_I,
 \qquad D_k(u_1\wedge u_I)=0.
\]
There are four pairs.

Choose grading shifts and cotangent partners that turn each
pair into the quartet just constructed. This specifies the
free paired sector. Give its quartets positive frequencies
$\omega_k$ with $c|k|\leq\omega_k\leq C|k|$ outside a finite
set, and define, for $\ell>0$,
\[
 W_k=\exp\!\left(\frac{\hbar C_k}{2\omega_k}\right),
 \qquad r_k=e^{-\ell\omega_k}.
\]

Here $C_k$ is the quartet contraction of
\eqref{mht:eq:quartet-wick}. The paired covariance is
$\hbar/(2\omega_k)$ and commutes with the mode differential.
Tensor products define preparation and transfer on polynomials
in finitely many modes. Smoothly smeared contractions converge
by rapid decay of their Fourier coefficients.

Complete the occupation basis with the vacuum as reference
vector. The logarithm of the absolute transfer trace is a
fixed finite multiple of
\[
 \sum_{k\ne0}2\log\frac{1+r_k}{1-r_k}<\infty.
\]
Indeed, $\sum_{k\ne0}e^{-\ell\omega_k}<\infty$ by lattice
shell counting. Absolute convergence permits the graded
product. Each quartet contributes one, so the nonzero-mode
supertrace is one.

At $k=0$ the differential vanishes. Before field shifts and
cotangent partners, the residual complex is
\[
 \Lambda^\bullet\langle ds,d\bar z_1,d\bar z_2\rangle,
\]
with eight components per coefficient copy. Its polarization
and state remain separate data. This free mixed Fourier
preparation has not identified those data with the compact
residual state of \eqref{mht:eq:residual-cylinder}.

\section{A scalar collar with infinitely many modes}

We next construct a positive scalar preparation on the compact
five-dimensional space. Its reference measure and vacuum energy
are fixed before the mode cutoff is removed. This makes it
possible to prove sewing and then add a nonnegative interaction.
The scalar sector is independently defined from the gauge
quartets and their residual cohomology.

Let $\mathbb T^5=(\mathbb R/2\pi\mathbb Z)^5$ with its flat
volume measure. Fix $m>0$ and numerical $\hbar>0$. Put
\[
 \Omega_0=\sqrt{-\Delta_{\mathbb T^5}+m^2},
 \qquad \Omega_0e_j=\omega_je_j,
\]
where $e_j$ is a real orthonormal Fourier basis.
Write $k_j$ for a frequency of $e_j$.
The eigenvalue multiset is
$\omega_k=\sqrt{|k|^2+m^2}$, $k\in\mathbb Z^5$.
The sine and cosine multiplicities give this same count.

Take independent centered Gaussian coordinates $x_j$ with
variances $c_j=\hbar/(2\omega_j)$.

For $s>2$, the random series $\phi=\sum_jx_je_j$ defines
a Gaussian probability measure $\mu$ on $H^{-s}(\mathbb T^5)$.
Indeed,
\begin{equation}\label{mht:eq:scalar-gaussian-domain}
 \mathbb E\|\phi\|_{H^{-s}}^2
   =\frac\hbar2\sum_{k\in\mathbb Z^5}
       \frac{(1+|k|^2)^{-s}}{\sqrt{|k|^2+m^2}}<\infty.
\end{equation}
Shell counting gives summands of order $n^{3-2s}$ after
grouping lattice points of radius $n$. The finite sums converge
in mean square in this Hilbert space. For real smooth $f$,
\begin{equation}\label{mht:eq:scalar-characteristic-functional}
 \int e^{i\langle\phi,f\rangle}\,d\mu(\phi)
 =\exp\!\left(-\frac\hbar4
                       \langle f,\Omega_0^{-1}f\rangle\right).
\end{equation}

\begin{proposition}[The exact Sobolev threshold]
\label{mht:prop:scalar-sobolev-threshold}
The measure $\mu$ is concentrated on $H^{-s}(\mathbb T^5)$
if and only if $s>2$.
\end{proposition}

\begin{proof}
The finite expectation in
\eqref{mht:eq:scalar-gaussian-domain} proves sufficiency.
For necessity put $a_j=(1+|k_j|^2)^{-s}c_j$ and
$\xi_j=x_j/\sqrt{c_j}$. Independence gives
\[
 \mathbb E\exp\!\left(-\sum_{j\leq K}a_j\xi_j^2\right)
       =\prod_{j\leq K}(1+2a_j)^{-1/2}.
\]

For $s\leq2$, shell counting gives $\sum_ja_j=\infty$.
Hence $\sum_j\log(1+2a_j)=\infty$. Indeed, either infinitely
many $a_j$ are bounded below, or the logarithms are comparable
to $a_j$ on the tail. The product tends to zero.

Bounded convergence shows that the exponential of the negative
full square sum is zero almost surely. Thus that square sum
is infinite almost surely, proving necessity.
\end{proof}

\subsection{The measure in the Mehler kernel}

For a collar of length $\ell>0$, let $r_j=e^{-\ell\omega_j}$.
Relative to the fixed measure $\mu$, the first $K$ modes have
kernel
\begin{equation}\label{mht:eq:scalar-finite-collar}
 p_\ell^{(K)}(x,y)=\prod_{j=1}^K(1-r_j^2)^{-1/2}
 \exp\!\left[-\sum_{j=1}^K
  \frac{r_j^2(x_j^2+y_j^2)-2r_jx_jy_j}
                              {2c_j(1-r_j^2)}\right].
\end{equation}
Unused modes have kernel one on the full probability space.
In one coordinate, $p_\ell(x,y)d\mu(y)$ is the Gaussian law
with mean $rx$ and variance $c(1-r^2)$.

The ordinary oscillator amplitude relative to Lebesgue measure is
\begin{equation}\label{mht:eq:scalar-oscillator-amplitude}
 Z_\ell(a,b)=
 \left(\frac{\omega}{2\pi\hbar\sinh(\omega\ell)}\right)^{1/2}
 \exp\!\left[-\frac{\omega((a^2+b^2)\cosh(\omega\ell)-2ab)}
                                  {2\hbar\sinh(\omega\ell)}\right].
\end{equation}
If $\psi_0^2$ is the Gaussian density of variance
$\hbar/(2\omega)$, direct substitution gives
\begin{equation}\label{mht:eq:scalar-ground-transform}
 p_\ell(a,b)=e^{\ell\omega/2}
                  \frac{Z_\ell(a,b)}{\psi_0(a)\psi_0(b)}.
\end{equation}
Thus \eqref{mht:eq:scalar-finite-collar} includes a change of
reference measure and a ground-energy subtraction.

\begin{theorem}\label{mht:thm:scalar-infinite-sewing}
The kernels $p_\ell^{(K)}$ converge in $L^2(\mu\otimes\mu)$
to a kernel $p_\ell$ that is positive almost everywhere.
For each pair of positive lengths, almost everywhere,
\begin{align}
 \int p_\ell(x,y)\,d\mu(y)&=1,\notag\\
 \int p_{\ell_1}(x,z)p_{\ell_2}(z,y)\,d\mu(z)
                         &=p_{\ell_1+\ell_2}(x,y).
       \label{mht:eq:scalar-infinite-sewing}
\end{align}
The associated operators $P_\ell$ form a strongly continuous,
self-adjoint Markov semigroup on $L^2(\mu)$.
At every positive length they are trace class, with
\begin{align}
 \operatorname{Tr}P_\ell
       &=\det(1-e^{-\ell\Omega_0})^{-1},\notag\\
 \|p_\ell\|_{L^2(\mu\otimes\mu)}^2
       &=\det(1-e^{-2\ell\Omega_0})^{-1}.
       \label{mht:eq:scalar-fredholm-determinants}
\end{align}
These Fredholm determinants converge on the one-particle
space $L^2(\mathbb T^5)$.
\end{theorem}

\begin{proof}
Composition of one-mode transitions gives mean $r_1r_2x$
and variance $c(1-r_1^2r_2^2)$. This proves finite-mode
sewing with its normalization. Direct Gaussian integration gives
\[
 \int p_r(x,y)^2\,d\mu(x)d\mu(y)=(1-r^2)^{-1}.
\]
Each independent coordinate factor has expectation one.
Consequently $p_\ell^{(K)}$ is a nonnegative martingale in
the first $K$ pairs of coordinates, with squared norm
$\prod_{j\leq K}(1-r_j^2)^{-1}$.

For every $a>0$, shell counting gives
$\sum_je^{-a\omega_j}<\infty$. Also $r_j\leq e^{-\ell m}<1$.
The squared norms are bounded, so the martingale converges
in $L^2$ to a nonnegative kernel.

To prove positivity, take the logarithm of the product.
The diagonal quadratic terms have summable absolute expectations.
The independent centered cross terms have summable variances,
and the constant logarithmic terms converge.
The logarithm is therefore finite almost everywhere.

Integration in one variable is an $L^2$ contraction, so
normalization passes to the limit. Kernel convergence gives
Hilbert--Schmidt convergence of the operators.
The finite operators are contractions. Their composition
identity therefore passes to the limit, proving
\eqref{mht:eq:scalar-infinite-sewing}.
Symmetry follows from the kernel formula.

Normalized Hermite polynomials in one coordinate have eigenvalues
$r^n$, as follows from their generating function.
Finite products give an orthonormal basis of $L^2(\mu)$:
conditional expectations onto finitely many coordinates
converge in $L^2$, and the finite-dimensional Hermite bases
are complete. The eigenvalues are
$\prod_jr_j^{n_j}$ for finite-support occupation sequences.
They tend to one as $\ell\downarrow0$ on each basis vector.
Density and contraction prove strong continuity.

Monotone summation of these positive eigenvalues gives
\[
 \operatorname{Tr}P_\ell
       =\prod_j\sum_{n\geq0}r_j^n
       =\prod_j(1-r_j)^{-1}<\infty.
\]
The one-particle operator $e^{-\ell\Omega_0}$ is trace class,
so this is its stated Fredholm determinant.
The identity $P_{2\ell}=P_\ell^*P_\ell$ gives the second one.
\end{proof}

For a Fourier cutoff $\|k\|_\infty\leq K$, let
$Z_\ell^{(K)}$ be the truncated trace and
$Z_\ell=\operatorname{Tr}P_\ell$. The shell cardinality is
\[
 (2n+1)^5-(2n-1)^5=160n^4+80n^2+2.
\]
Since $\omega_k\geq\|k\|_\infty$ and
$-\log(1-x)\leq x/(1-x)$,
\begin{equation}\label{mht:eq:scalar-fredholm-tail}
 0\leq\log Z_\ell-\log Z_\ell^{(K)}
 \leq\frac{\sum_{n>K}(160n^4+80n^2+2)e^{-\ell n}}
                                  {1-e^{-\ell(K+1)}}.
\end{equation}
The massive zero mode is included. At $m=0$, its variance
and trace factor are undefined in this construction.
It must then be prepared as a residual variable.

\section{An interaction at fixed positive smoothing}

Fix $\sigma>0$. Define a smoothing operator and potential by
\begin{equation}\label{mht:eq:scalar-smoothed-potential}
 R_\sigma=e^{-\sigma\Omega_0},\qquad
 V_\sigma(\phi)=\frac1{24}\int_{\mathbb T^5}
                                      (R_\sigma\phi)(z)^4\,dz.
\end{equation}
The map $R_\sigma$ sends $H^{-s}$ continuously into every
fixed smooth Sobolev space. Thus $V_\sigma$ is continuous,
finite and nonnegative on $H^{-s}$. Smoothing bounds it by
a constant times $\|\phi\|_{H^{-s}}^4$.
All its Gaussian moments are finite.

For paths fix $s>2$. Let $\Phi_t$ be the stationary
Ornstein--Uhlenbeck process with coordinate covariance
\[
 \mathbb E(\Phi_j(t)\Phi_j(t'))
                  =c_je^{-\omega_j|t-t'|}.
\]

\begin{proposition}\label{mht:prop:scalar-time-regularity}
The process has a continuous $H^{-s}(\mathbb T^5)$ modification,
locally H\"older of every order $\alpha<\min\{1/2,s-2\}$.
For $0<h\leq1/2$ its increment variance satisfies
\begin{equation}\label{mht:eq:scalar-time-increments}
 \mathbb E\|\Phi_{t+h}-\Phi_t\|_{H^{-s}}^2\leq C_s
 \begin{cases}
 h^{2s-4},&2<s<5/2,\\
 h\log(e/h),&s=5/2,\\
 h,&s>5/2.
 \end{cases}
\end{equation}
\end{proposition}

\begin{proof}
The increment variance is
\[
 2\sum_j(1+|k_j|^2)^{-s}c_j(1-e^{-\omega_jh}).
\]
The bound $1-e^{-x}\leq\min\{1,x\}$ and shell counting reduce it
to a constant times $\sum_{n\geq1}n^{3-2s}\min\{1,nh\}$,
plus a bounded multiple of $h$ from the massive zero mode.
Splitting this sum at $n=h^{-1}$ gives
\eqref{mht:eq:scalar-time-increments}.
For $s>5/2$, the simpler bound
$\hbar h\sum_j(1+|k_j|^2)^{-s}$ also converges.

For a centered Hilbert-valued Gaussian $G$, diagonalize its
covariance and apply the $L^q$ triangle inequality to its squared
norm. This gives
$\mathbb E\|G\|^{2q}\leq C_q(\mathbb E\|G\|^2)^q$.
Apply this to the increments and use the moment continuity
criterion with arbitrarily large integers $q$.
At $s=5/2$, the logarithm is bounded by every negative power
of $h$, giving each strict exponent $\alpha<1/2$.
\end{proof}

The transition kernel of this continuous modification is $p_t$.

The continuous path space is separable, so regular conditional
Gaussian bridge laws exist given $\Phi_0=x$, $\Phi_\ell=y$.
Their endpoint law is $p_\ell(x,y)d\mu(x)d\mu(y)$.
Positivity of $p_\ell$ makes this measure equivalent to
$\mu\otimes\mu$. For $g\geq0$ define
\begin{equation}\label{mht:eq:scalar-interacting-kernel}
 p_\ell^{g,\sigma}(x,y)=p_\ell(x,y)
  \mathbb E_{x,y}^{\mathrm{bridge}}
       \exp\!\left(-g\int_0^\ell V_\sigma(\Phi_t)\,dt\right).
\end{equation}
Here $g$ is the coefficient in the Euclidean weight.
An action coupling $g_{\mathrm{phys}}$ corresponds to
$g=g_{\mathrm{phys}}/\hbar$.

\begin{theorem}\label{mht:thm:scalar-interacting-sewing}
For $g\geq0$ and $\sigma>0$, the kernel
\eqref{mht:eq:scalar-interacting-kernel} is symmetric, belongs
to $L^2(\mu\otimes\mu)$, and satisfies
$0\leq p_\ell^{g,\sigma}\leq p_\ell$.
For each pair of positive lengths, almost everywhere,
\begin{equation}\label{mht:eq:scalar-interacting-sewing}
 \int p_{\ell_1}^{g,\sigma}(x,z)
       p_{\ell_2}^{g,\sigma}(z,y)\,d\mu(z)
                         =p_{\ell_1+\ell_2}^{g,\sigma}(x,y).
\end{equation}
The operators form a strongly continuous, self-adjoint,
positivity-preserving contraction semigroup.
At positive length they are trace class and satisfy
\begin{equation}\label{mht:eq:scalar-interacting-trace}
 \operatorname{Tr}P_\ell^{g,\sigma}
   =\|P_{\ell/2}^{g,\sigma}\|_{\mathrm{HS}}^2
   \leq\det(1-e^{-\ell\Omega_0})^{-1}.
\end{equation}
\end{theorem}

\begin{proof}
First use a bounded nonnegative potential. Its path weight
lies between zero and one. Conditioning at an intermediate
time and using the Markov property gives the sewing law.
Reversal of the stationary path gives symmetry.

Apply this construction to $V_\sigma\wedge M$ and let
$M\to\infty$. On every continuous path, the potential
integral is finite by continuity and compactness of the time
interval. The weights decrease to the weight in
\eqref{mht:eq:scalar-interacting-kernel}.
Domination by $p_\ell$ gives convergence in $L^2$.
In the sewing integral the product is dominated by the free
product, whose integral is $p_{\ell_1+\ell_2}$.
Dominated convergence proves the claimed sewing law.

The same kernel is obtained by spatial mode cutoffs.
Fourier projections converge uniformly on the compact image
of a continuous $H^{-s}$ path. After smoothing, the potentials
and their integrated weights therefore converge.
The finite-coordinate marginal of the full bridge is precisely
the corresponding finite-mode bridge. Its conditional weights
converge almost everywhere and are bounded by one.
Together with $p_\ell^{(K)}\to p_\ell$ in $L^2$, this proves
the cutoff assertion.

Kernel domination gives the Hilbert--Schmidt bound.
Symmetry and sewing imply
\[
 P_\ell^{g,\sigma}
       =(P_{\ell/2}^{g,\sigma})^*P_{\ell/2}^{g,\sigma}.
\]
This proves trace class, operator positivity and
\eqref{mht:eq:scalar-interacting-trace}.
The inequality $P_\ell^{g,\sigma}1\leq1$ and symmetry give
contraction on $L^2(\mu)$.
On bounded continuous cylinder functions, path continuity
and bounded convergence prove strong continuity at zero.
Density and contraction extend it to all of $L^2(\mu)$.
\end{proof}

\begin{corollary}\label{mht:cor:scalar-trace-convergence}
At fixed $g\geq0$, $\sigma>0$ and $\ell>0$, the interacting
mode-cutoff operators converge to $P_\ell^{g,\sigma}$ in trace norm.
The finite kernels are extended by the constant kernel one in
the unused modes.
\end{corollary}

\begin{proof}
The proof of Theorem~\ref{mht:thm:scalar-interacting-sewing}
gives Hilbert--Schmidt convergence at length $\ell/2$.
Write $A_K$ and $A$ for the half-length operators.
The finite sewing identity includes the Gaussian
conditional-expectation projection on unused modes.
Consequently its length-$\ell$ operator is $A_K^2$,
and the limiting sewing identity gives $A^2$.

Decompose the difference as
$(A_K-A)A_K+A(A_K-A)$.
The trace-ideal product estimate gives
\[
 \|A_K^2-A^2\|_1
 \leq\|A_K-A\|_2(\|A_K\|_2+\|A\|_2),
\]
which tends to zero.
\end{proof}

At positive coupling these operators are generally sub-Markov:
their mass can be less than one. Dividing each interval kernel
by its own mass would generally change the sewing law.
The path weights in \eqref{mht:eq:scalar-interacting-kernel}
remain part of the amplitude.

\subsection{The first insertion and the ordering of composites}

Conditioning the stationary covariance on both endpoints gives
\begin{align}
 m_j(t)&=\frac{\sinh(\omega_j(\ell-t))}{\sinh(\omega_j\ell)}x_j
           +\frac{\sinh(\omega_jt)}{\sinh(\omega_j\ell)}y_j,
       \label{mht:eq:scalar-bridge-mean}\\
 v_j(t)&=\frac{\hbar\sinh(\omega_jt)
                            \sinh(\omega_j(\ell-t))}
                            {\omega_j\sinh(\omega_j\ell)}.
       \label{mht:eq:scalar-bridge-variance}
\end{align}
For a direct derivation, put $r=e^{-\omega_jt}$ and
$s=e^{-\omega_j(\ell-t)}$. The endpoint covariance is
$c_j\left(\begin{smallmatrix}1&rs\\rs&1\end{smallmatrix}\right)$
and the cross covariance is $c_j(r,s)$.
Multiplying by the inverse endpoint matrix gives the mean.
The Schur complement is
$c_j(1-r^2)(1-s^2)/(1-r^2s^2)$, giving the variance.

Define the smoothed mean and pointwise variance by
\[
 m_\sigma(t,z)=\sum_je^{-\sigma\omega_j}m_j(t)e_j(z),
 \qquad
 v_\sigma(t,z)=\sum_je^{-2\sigma\omega_j}v_j(t)e_j(z)^2.
\]
These series converge with every fixed spatial derivative.
The conditional variance is bounded by $c_j$, and the
smoothed mean is bounded by the endpoint $H^{-s}$ norms.
Consequently all moments of the integrated potential are
finite for almost every endpoint pair.

\begin{proposition}\label{mht:prop:scalar-first-insertion}
The first right response at zero coupling is
\begin{align}
 -\left.\partial_g\log
       \frac{p_\ell^{g,\sigma}(x,y)}{p_\ell(x,y)}\right|_{g=0+}
  =\frac1{24}\int_0^\ell dt\int_{\mathbb T^5}dz\,
       \bigl(m_\sigma^4+6m_\sigma^2v_\sigma+3v_\sigma^2\bigr).
       \label{mht:eq:scalar-first-insertion}
\end{align}
All higher right derivatives at zero exist as finite bridge
moments and cumulants.
\end{proposition}

\begin{proof}
Put $\mathcal V=\int_0^\ell V_\sigma(\Phi_t)dt$.
The quotient $(1-e^{-g\mathcal V})/g$ is bounded by
$\mathcal V$ and tends to it as $g\downarrow0$.
Dominated convergence gives the derivative of its conditional
expectation. At zero coupling that expectation is one, so
the negative logarithmic derivative is
$\mathbb E_{x,y}^{\mathrm{bridge}}\mathcal V$.

For a Gaussian of mean $m$ and variance $v$,
$\mathbb E X^4=m^4+6m^2v+3v^2$.
Integrating this identity proves the formula.
Higher finite moments justify repeated right differentiation.
Differentiating the logarithm gives their cumulants.
\end{proof}

The last two terms retain fluctuations about the conditional
mean. A change of covariance has its own polynomial ordering
rule. For a scalar variance $C$, define
$:F:_C=e^{-C\partial_x^2/2}F$. Then
\begin{equation}\label{mht:eq:scalar-wick-ordering}
 :F:_{C_2}=e^{-(C_2-C_1)\partial_x^2/2}:F:_{C_1},
 \qquad :x^4:_C=x^4-6Cx^2+3C^2.
\end{equation}
The exponentials terminate on polynomials.

On smeared kernels the same identity uses the corresponding
covariance contraction. It does not define unsmoothed
coincident products.

\section{Divergence and Borel reconstruction}

The finite bridge moments determine an asymptotic expansion at
zero coupling. Their factorial growth also determines a Borel
reconstruction of the positive-coupling kernel. Write
$X=\mathbb T^5$, $\mathcal V=\int_0^\ell V_\sigma(\Phi_t)\,dt$ and
$F_{\ell,x,y}(g)=\mathbb E_{x,y}^{\mathrm{bridge}}e^{-g\mathcal V}$.
This expectation is defined for $\operatorname{Re}g\geq0$.
Let $J_0(z)=\sum_{n\geq0}(-1)^n(z^2/4)^n/(n!)^2$ be the
Bessel function of order zero.

\begin{lemma}\label{mht:lem:scalar-bessel-reconstruction}
Let $\mathcal V\geq0$ be a random variable with
$\mathbb E e^{\varepsilon\sqrt{\mathcal V}}<\infty$ for some $\varepsilon>0$.
Put $M_n=\mathbb E \mathcal V^n$ and
$F(g)=\mathbb E e^{-g\mathcal V}$ for $g\geq0$.
The Borel germ of its formal expansion is
\[
 \widehat F(\zeta)=\sum_{n\geq0}
             \frac{(-1)^nM_n}{(n!)^2}\zeta^n.
\]
It continues along the nonnegative ray by
\begin{equation}\label{mht:eq:scalar-bessel-transform}
 \widehat F(\zeta)=\mathbb E J_0(2\sqrt{\zeta \mathcal V}),\qquad
 |\widehat F(\zeta)|\leq1\quad(\zeta\geq0),
\end{equation}
and for every $g>0$,
\begin{equation}\label{mht:eq:scalar-borel-laplace}
 F(g)=\int_0^\infty e^{-t}\widehat F(gt)dt.
\end{equation}
The continuation is analytic wherever
$2|\operatorname{Im}\sqrt\zeta|<\varepsilon$, together with its
germ at zero.
\end{lemma}
\begin{proof}
The function $J_0(2\sqrt{\zeta \mathcal V})$ is entire in $\zeta$, by its
power series. The integral representation
\[
 J_0(z)=\frac1{2\pi}\int_0^{2\pi}e^{iz\cos\theta}\,d\theta
\]
gives $|J_0(z)|\leq e^{|\operatorname{Im}z|}$.
For small $|\zeta|$, absolute summation is dominated by
$e^{2\sqrt{|\zeta|}\sqrt{\mathcal V}}$. This proves the germ identity.
The same bound on compact subsets proves the stated analytic
continuation, while for nonnegative $\zeta$ the bound is one.

For fixed $w\geq0$, termwise integration is justified by the
convergent sum $\sum_n(gw)^n/n!=e^{gw}$ and gives
\[
 \int_0^\infty e^{-t}J_0(2\sqrt{gtw})dt=e^{-gw}.
\]
To average this identity, use $|J_0(2\sqrt{gt\mathcal V})|\leq1$.
Fubini now requires only the integrable bound $e^{-t}$ and proves
\eqref{mht:eq:scalar-borel-laplace}.
No exponential moment $\mathbb E e^{g\mathcal V}$ is required.
\end{proof}

\begin{theorem}\label{mht:thm:scalar-quartic-borel}
Fix $m,\hbar,\ell,\sigma>0$ and $s>2$ in
\eqref{mht:eq:scalar-interacting-kernel}. For almost every endpoint pair $(x,y)$ there
are constants $0<c_-\leq c_+<\infty$ such that, for $n\geq1$,
\begin{equation}\label{mht:eq:scalar-gevrey}
 c_-^n(n!)^2\leq\mathbb E_{x,y}^{\mathrm{bridge}}\mathcal V^n\leq c_+^n(n!)^2.
\end{equation}
The coupling series of $F_{\ell,x,y}$ has radius zero and exact
Gevrey order one. Its coefficients are bounded by an exponential
times $n!$. No exponential multiple of $(n!)^a$, with $a<1$,
bounds them. It is an asymptotic expansion on $\operatorname{Re}g\geq0$,
with a remainder bounded by $c_+^N N!|g|^N$ after $N$ terms.

Its endpointwise Borel continuation and Laplace integral are given by
Lemma~\ref{mht:lem:scalar-bessel-reconstruction}. Multiplication by $p_\ell$ gives
$L^2(\mu\otimes\mu)$ kernels along the positive Borel ray, whose
Bochner Laplace integral is exactly $p_\ell^{g,\sigma}$.
Thus this reconstruction preserves the sewing of
Theorem~\ref{mht:thm:scalar-interacting-sewing}.
\end{theorem}
\begin{proof}
Write $D=\ell\operatorname{vol}(X)$ and
$u(t,z)=(R_\sigma\Phi_t)(z)$. Its conditional mean has a finite
uniform bound $M$, and its conditional pointwise variance has a
finite uniform bound $v_*$, because $\sigma>0$ and $x,y\in H^{-s}$.
For a standard real Gaussian $Z$, let $b_n=\mathbb E Z^{4n}=(4n-1)!!$.
The exact ratio
\[
 \frac{b_{n+1}}{b_n}=(4n+1)(4n+3)
\]
lies between $3(n+1)^2$ and $16(n+1)^2$. Hence
$3^n(n!)^2\leq b_n\leq16^n(n!)^2$.
Using $|a+b|^{4n}\leq2^{4n-1}(|a|^{4n}+|b|^{4n})$ and Jensen's
inequality on the finite measure space $[0,\ell]\times X$ yields
\[
 \mathbb E \mathcal V^n\leq
 \left[\frac{D}{24}\,16(M^4+16v_*^2)\right]^n(n!)^2.
\]

For a lower bound use the space-time average
$Y=D^{-1}\int_0^\ell\int_Xu(t,z)dzdt$.
H\"older's inequality gives $\mathcal V\geq(D/24)Y^4$.
The constant spatial mode is $e_0=\operatorname{vol}(X)^{-1/2}$, so
\[
 Y=\frac{e^{-\sigma m}}{\ell\sqrt{\operatorname{vol}(X)}}
                        \int_0^\ell\Phi_0(t)dt.
\]
Its conditional variance is
\begin{equation}\label{mht:eq:scalar-average-variance}
 v_Y=\frac{\hbar e^{-2\sigma m}}{\ell^2\operatorname{vol}(X)m^2}
       \left(\ell-\frac2m\tanh\frac{m\ell}{2}\right)>0.
\end{equation}

To obtain this formula, the covariance of the pinned constant mode
is $\hbar$ times the Dirichlet Green kernel of
$-\partial_t^2+m^2$. Its integral solves
$(-\partial_t^2+m^2)f=1$, $f(0)=f(\ell)=0$, and is
\[
 f(t)=\frac1{m^2}
       \left(1-\frac{\cosh(m(t-\ell/2))}{\cosh(m\ell/2)}\right).
\]
Integrating $f$ proves \eqref{mht:eq:scalar-average-variance}.
Positivity follows from $\tanh a<a$ for $a>0$.

All terms in the even Gaussian moment expansion are nonnegative,
so $\mathbb E Y^{4n}\geq v_Y^{2n}b_n$.
This proves the lower bound with $c_-=3Dv_Y^2/24$.

Dividing \eqref{mht:eq:scalar-gevrey} by $n!$ proves the radius and exact
factorial order. The integral Taylor remainder for $e^{-g\mathcal V}$ has
absolute value at most $|g\mathcal V|^N/N!$ on the closed right half-plane,
which proves the asymptotic estimate.

The upper moment bound also implies exponential square-root
integrability: use
\[
 e^{a\sqrt{\mathcal V}}\leq2\cosh(a\sqrt{\mathcal V}),\qquad
 \mathbb E\cosh(a\sqrt{\mathcal V})
 =\sum_{n\geq0}\frac{a^{2n}\mathbb E \mathcal V^n}{(2n)!}<\infty
\]
for $a^2c_+<1$, since $(2n)!\geq(n!)^2$.
Lemma~\ref{mht:lem:scalar-bessel-reconstruction} applies.

Along the positive ray the reconstructed kernel is bounded in
absolute value by $p_\ell$. Dominated convergence makes it continuous
in $L^2$. Its Bochner integral therefore exists
in $L^2$, and the endpointwise identity identifies it with
\eqref{mht:eq:scalar-interacting-kernel}. The resulting equality of kernels transfers
their actual convolution law.
\end{proof}

The analytic germ in this theorem is endpointwise. Its radius may
depend on the endpoints. The positive-ray Bochner integral uses the
uniform domination by $p_\ell$. No endpoint-independent analytic
radius in $L^2$ is needed for that assertion.

The second connected response has an explicit form as well.
Let $X_1,X_2$ be smoothed field values with means $a,b$,
variances $v,w$ and covariance $C$. Gaussian expansion gives
\begin{equation}\label{mht:eq:scalar-quartic-covariance}
 \begin{split}
 \operatorname{Cov}(X_1^4,X_2^4)={}&
 16ab(a^2+3v)(b^2+3w)C
 +72(a^2+v)(b^2+w)C^2\\
 &+96abC^3+24C^4.
 \end{split}
\end{equation}

For a proof, expand $X_1^4$ in the Wick polynomials of its centered
Gaussian part. The coefficients of Wick degrees $1,2,3,4$ are
$4a(a^2+3v)$, $6(a^2+v)$, $4a$, $1$. Cross expectation of equal
Wick degrees $r$ is $r!C^r$, while unequal degrees have zero cross
expectation. This proves \eqref{mht:eq:scalar-quartic-covariance}.
Its double integral divided by $24^2$ is
$\partial_g^2\log F_{\ell,x,y}(0+)=\operatorname{Var}(\mathcal V)$.

\section{A unitary comparison of the scalar states}

The scalar kernel has a Hilbert Fock realization with the same
interaction and vacuum normalization. Construct its Hamiltonian
on the symmetric Hilbert Fock space
\[
 \mathcal F_{\Omega_0}
   =\widehat{\bigoplus}_{n\geq0}\operatorname{Sym}^n\ell^2(\{j\}),
\]
where the summands are symmetric Hilbert tensor powers.
Its orthonormal occupation basis $|\mathbf n\rangle$ is indexed by
finitely supported sequences of nonnegative integers.
The creation and annihilation operators $b_j^\dagger,b_j$ obey
$[b_i,b_j^\dagger]=\delta_{ij}$, with all other commutators zero.
The energy $d\Gamma(\Omega_0)$ has eigenvalue $\sum_j\omega_jn_j$.

For $\sigma>0$ define the selfadjoint Segal field
\[
 Q_\sigma(z)=\sum_j e^{-\sigma\omega_j}e_j(z)
                           \sqrt{c_j}(b_j+b_j^\dagger).
\]
The real coefficient sequence is square-summable uniformly in $z$.
On $\mathcal F_{\Omega_0}$ define the maximal operator
\[
 \mathsf T_\sigma:\mathcal F_{\Omega_0}\longrightarrow L^2(X;\mathcal F_{\Omega_0}),
 \qquad (\mathsf T_\sigma\psi)(z)=Q_\sigma(z)^2\psi,
\]
whose domain consists of vectors for which the displayed expression
is defined almost everywhere and is square-integrable.
This operator is closed. A convergent graph sequence has a
subsequence converging pointwise almost everywhere in the target.
Closedness of each $Q_\sigma(z)^2$ identifies the limit.

On particle sectors of number at most $K$, the creation and
annihilation bounds give
$\|Q_\sigma(z)^2\psi\|\leq C_\sigma(K+2)\|\psi\|$ uniformly in $z$.
Thus finite occupation vectors lie in its domain.

For $g>0$ define the quadratic form
\begin{equation}\label{mht:eq:scalar-fock-form}
 \begin{aligned}
 q_{g,\sigma}(\psi)
   &=\|(d\Gamma(\Omega_0))^{1/2}\psi\|^2
       +\frac g{24}\|\mathsf T_\sigma\psi\|^2,\\
 D(q_{g,\sigma})
   &=D((d\Gamma(\Omega_0))^{1/2})\cap D(\mathsf T_\sigma).
 \end{aligned}
\end{equation}
It is nonnegative and closed: its form norm is the joint graph
norm of two closed operators. Its domain contains the algebraic
span of the occupation basis and is therefore dense.

Let $K_{g,\sigma}$ be its
associated nonnegative selfadjoint operator. For $g=0$ set
$K_{0,\sigma}=d\Gamma(\Omega_0)$ on $D(d\Gamma(\Omega_0))$.
The form representation and increasing-form convergence apply
to these densely defined nonnegative forms.\footnote{See
\href{https://webhomes.maths.ed.ac.uk/~v1ranick/papers/kato1.pdf}
{T.~Kato, \emph{Perturbation Theory for Linear Operators}}, second
edition, corrected printing. The results are VI.2.1, VI.2.6,
VI.2.23 and VIII.3.13a.}

Let $\mathrm{He}_n$ be the probabilists' Hermite polynomial, defined by
$e^{ty-t^2/2}=\sum_n\mathrm{He}_n(y)t^n/n!$. For a finitely supported
occupation sequence set
\[
 h_{\mathbf n}(x)=\prod_j
       \frac{\mathrm{He}_{n_j}(x_j/\sqrt{c_j})}{\sqrt{n_j!}}.
\]

\begin{proposition}\label{mht:prop:scalar-fock-unitary}
The map
\[
 \mathcal U:L^2(\mu)\longrightarrow\mathcal F_{\Omega_0},
 \qquad \mathcal U h_{\mathbf n}=|\mathbf n\rangle
\]
extends to a unitary isomorphism. On polynomial cylinder vectors it
preserves the following operators and their common commutation relations:
\begin{align*}
 \mathcal U1&=|0\rangle,\\
 \mathcal U\,x_j\,\mathcal U^{-1}
       &=\sqrt{c_j}(b_j+b_j^\dagger),\\
 \mathcal U\,\sqrt{c_j}\partial_{x_j}\,\mathcal U^{-1}&=b_j,\\
 \mathcal UP_\ell\mathcal U^{-1}&=e^{-\ell d\Gamma(\Omega_0)}.
\end{align*}
At fixed smoothing, the multiplication form of $V_\sigma$ is carried to
$\|\mathsf T_\sigma\psi\|^2/24$, with equality of the full form
domains. For $g\geq0$ one has
\begin{equation}\label{mht:eq:scalar-interacting-unitary}
 \mathcal U P_\ell^{g,\sigma}\mathcal U^{-1}
                       =e^{-\ell K_{g,\sigma}}.
\end{equation}
This comparison preserves operator composition, adjoints, traces
and normalized positive states.
\end{proposition}
\begin{proof}
The generating function and Gaussian integration give
$\mathbb E\mathrm{He}_n(Z)\mathrm{He}_r(Z)=n!\delta_{nr}$.
Finite products form an orthonormal basis because conditional
expectations onto finitely many coordinates converge in $L^2$ and
finite-dimensional Hermite polynomials are complete. Thus $\mathcal U$
is unitary. The recurrence
$y\mathrm{He}_n=\mathrm{He}_{n+1}+n\mathrm{He}_{n-1}$ and
$\mathrm{He}'_n=n\mathrm{He}_{n-1}$ prove the field and derivative
identities. In particular, the Gaussian adjoint of
$\sqrt{c_j}\partial_{x_j}$ is
$x_j/\sqrt{c_j}-\sqrt{c_j}\partial_{x_j}$, matching $b_j^\dagger$.
The free heat identity follows on the same basis and hence on the
Hilbert completion.

Exponential smoothing makes the coefficient sequence of
$Q_\sigma(z)$ square-summable, with all fixed spatial derivatives.
The multiplication operator $R_\sigma\phi(z)$ is its unitary image
under the preceding identities, first on cylinders and then by closure.

Finite occupation vectors are analytic for every real Segal field,
by the creation and annihilation bounds. Their Gaussian polynomials
are analytic for the corresponding real linear multiplication
operator, so equality on this core identifies the selfadjoint closures.

The spectral theorem identifies the full domain of $Q_\sigma(z)^2$
with multiplication by $(R_\sigma\phi(z))^2$. Tonelli's theorem then
identifies the domain of $\mathsf T_\sigma$ with $D(V_\sigma^{1/2})$.

Let $H_{\mathrm{OU}}$ be the nonnegative generator of $P_\ell$.
Together with the free heat identity, this identifies the full
closed form \eqref{mht:eq:scalar-fock-form} with the Gaussian form
\[
 \|H_{\mathrm{OU}}^{1/2}f\|^2
                       +g\|V_\sigma^{1/2}f\|^2.
\]

For each bounded potential $\min(V_\sigma,M)$, expansion in its
bounded multiplication operator and the Markov law give the
Feynman--Kac semigroup. The forms increase to the displayed closed
form as $M\to\infty$. Their associated resolvents and semigroups
converge strongly, while their kernels converge by bounded
convergence of the path weights. This is precisely the preparation
in Theorem~\ref{mht:thm:scalar-interacting-sewing}, proving
\eqref{mht:eq:scalar-interacting-unitary}.

Trace-class positivity and unitary
invariance of the trace then give, for every bounded $A$,
\[
 \frac{\operatorname{Tr}(P_\ell^{g,\sigma}A)}{\operatorname{Tr} P_\ell^{g,\sigma}}
 =\frac{\operatorname{Tr}\bigl(e^{-\ell K_{g,\sigma}}
                   (\mathcal UA\mathcal U^{-1})\bigr)}
        {\operatorname{Tr} e^{-\ell K_{g,\sigma}}}.
\]
This is an equality of prepared states under the specified unitary,
with the same reference vacuum and coupling.
\end{proof}

\section{The comparison with native quantum fields}

The compact residual state, determinant-line Gaussian,
BRST quartet and positive scalar collar each have a specified
domain and sewing law. Their constructions supply concrete
targets for the native comparison of
Problem~\ref{mht:qn:native-quantum-elimination}.
That comparison must preserve more than the vacuum amplitude.

The invariant functions $\mathcal O_r$ of
Chapter~\ref{mht:ch:nonlinear-neumann} already lift classical
matrix invariants to smooth Neumann fields.
When the chosen real quadratic form is
$\langle J,J\rangle=c_G\operatorname{tr}(J^2)$, its potential
is $vc_G\mathcal O_2$. Its quantum image must specify the
ordering, including the lower-Casimir corrections.
The half-density computation also requires a map between
the holonomy coordinate $X$ and the residual field $J$.

\begin{problem}\label{mht:qn:native-prepared-state}
Construct a renormalized observable complex for the native
Neumann fields and a boundary state map from its action.
Specify a common regulator, the determinant or Pfaffian line,
the residual measure and the boundary BRST polarization.
Compare the resulting invariant degree-zero states with
$L^2(G,dg)^G$ and the heat operator
\eqref{mht:eq:prepared-heat-time}.

For a cut surface $U=U_1\cup_\Sigma U_2$, let $F_i$ be
regular observables supported inside $U_i$.
Construct source-dependent sewing of the form
\[
 \Pi_U(F_1F_2)=\operatorname{BV}_{\mathrm{res}}
       \bigl\langle\Pi_{U_1}(F_1),\Pi_{U_2}(F_2)\bigr\rangle_\Sigma,
\]
where the pairing includes line transport and the residual
BV integral satisfies Stokes. Require adjoint boundary
differentials and compatibility with the supported insertion
operations. Determine the disk normalization, vacuum energy,
holonomy density and ordering of the lifted $\mathcal O_r$.
If the prepared states differ, compute their first invariant
defect modulo the allowed local cyclic counterterms.
\end{problem}

The scalar interaction poses a separate analytic limit:
construct local counterterms and a state for which its sewing
maps converge as $\sigma\downarrow0$.
At fixed positive smoothing, the divergent coupling series has
a Borel reconstruction and the scalar Fock comparison is unitary.
The local limit and the comparison with the physical BRST fields,
their residual sector and their boundary pairing remain to be
constructed.
